\documentclass{article}
\usepackage{pathfinder-readable}

\title{When a Wrist Frame Needs a New Observation}

\begin{document}
\maketitle

\begin{abstract}
Paper P studies a donor wrist as a handover frame when the held object is hidden, while paper Q surveys
how predictions should be judged by available information, task error and cost. The thread asked whether
the receiver's full visible history can still leave different object poses that call for incompatible
grasps. The peers derived conditional decision and sensing bounds, but found no measured handover
evidence for the premise. The verifier paused the thread because P already gives the basic obstruction
and the added calculations do not yet support a distinct paper section. A matched-history experiment with
receiving outcomes and sensing costs is needed to resume it.
\end{abstract}

\section{Introduction}

Paper P is a local manuscript about handover when the object is occluded. It describes using the donor
wrist pose as a substitute task frame. It also gives a conditional obstruction: if the receiver sees the
same wrist, motor-encoder and action history in two trials, but the hidden object orientations need
incompatible approaches, the receiver cannot choose correctly in both. P does not report the matched
histories or receiving outcomes needed to show that this case occurs in practice \cite{P}.

Paper Q surveys learning-augmented algorithms, which use predictions while accounting for their errors.
It says a deployed predictor can use only information available when it predicts. It also asks for an
error measure tied to the task and for the cost of obtaining a prediction to be counted. Q is a survey,
not a handover study \cite{Q}.

The connection was a decision question. In-hand object rotation can change the object's pose relative to
a fixed wrist. The peers asked whether the receiver's complete visible history then leaves a choice
between object poses for which no single grasp works. They used Q's rules for observable predictions,
task loss and query cost to make P's handover question more precise.

\section{What was done}

The peers first proposed treating the hidden object pose as a prediction target. They then moved from
pose error to receiving success. A difference in pose matters only if it changes which actions can
succeed. They wrote a common-grasp condition and a conditional failure risk for a receiver that uses the
full visible history. This followed P's matched-observation argument but also covered cases in which a
robust action succeeds across the possible poses.

They next checked a prediction-set certificate. If a covered set of poses admits one action that succeeds
throughout the set, failure can occur only when the true pose falls outside it. One peer also considered
a bound on grasp-point position error. The other pointed out that rotation about that point can leave its
position unchanged while changing the required receiving orientation. They retained the full-pose,
task-specific condition.

The peers then corrected two proposed transfers from Q. Q's formal consistency claim still makes sense
for either hidden state supplied with its own hypothetical exact prediction. The difficulty is deploying
one predictor on identical observations in both states. They also found that a multiplicative cost ratio
is ill suited to a one-shot handover with binary failure: a receiver told the true pose may have zero
failure cost. They therefore used additive failure and sensing costs. Finally, they checked related work
and sharpened the proposed physical test.

\section{Results}

Let \(Z=z\) be everything visible to the receiver at its decision, including wrist pose, encoder readings
and action history. Let \(G\) be the hidden object pose relative to the wrist, and let \(S(G)\) be the
set of receiving actions that succeed at that pose. For binary failure, the infimum risk over actions
using only \(z\) is
\[
R^*(z)=1-\sup_a\Pr\{a\in S(G)\mid Z=z\}.
\]
This is the lowest failure chance that such actions can attain or approach.
The peers derived this form in the ledger and checked its two-state case. It expresses the issue in terms
of receiving decisions, rather than pose spread alone.

For a guarantee across all poses compatible with \(z\), a successful history-only action exists exactly
when their sets \(S(G)\) have a common member. P already states the obstruction for two matched states
that require incompatible approaches \cite{P}. The positive condition explains why pose variation alone
need not cause failure: a common safe grasp would settle the choice.

The peers also checked a limited prediction-set result. Suppose a pose set computed from \(Z\) contains
the true \(G\) with marginal probability at least \(1-\alpha\), where \(\alpha\) is the allowed miss
probability. If the receiver chooses an action that succeeds for every pose in that set, and the object
remains fixed until the action, its marginal failure probability is at most \(\alpha\). This follows
because failure implies that the true pose was outside the set. It says nothing about safety for each
particular history.

For a simple sensing calculation, the peers assumed two equally likely, fixed hidden poses with nonempty
but disjoint successful-action sets. They called the state-dependent sensor reading \(W\), and its two
conditional laws \(K_0\) and \(K_1\). Without that reading, the minimum failure probability is \(1/2\).
With it, the minimum is \((1-\operatorname{TV}(K_0,K_1))/2\), where \(\operatorname{TV}\) is the total
variation distance between those laws. If failure costs \(M\) and sensing costs \(\kappa\), with no other
costs, sensing reduces expected additive cost exactly when \(\kappa<M\operatorname{TV}(K_0,K_1)/2\). The
peers checked this binary decision calculation independently. Repeating a predictor on the same history
with randomness independent of the hidden state supplies no new state information; the reading must
depend on the hidden state. This is a conditional
calculation, not a measured handover result or a direct use of Q's predictor-call ratio \cite{Q}.

\section{What did not hold}

An early suggestion said Q's consistency condition was vacuous on matched states. The peers withdrew that
wording: Q tests consistency on an instance paired with an exact prediction, while the handover limit
concerns whether one deployed predictor can produce distinct exact poses from the same input. They also
rejected grasp-point position error alone as a certificate for an orientation-sensitive grasp. The
one-shot failure model required additive accounting because its clairvoyant optimum can be zero.

The physical premise remained unsettled. P supplies no measured distribution of object pose given a
complete wrist, encoder and action history. Approximate matches among continuous histories would need a
justified tolerance and would not prove exact indistinguishability. A common safe grasp could remove the
action conflict. A sensing action that moves the object would require a model of that change before the
fixed-state calculation could be used.

The peers' prior-work checks also limited the claim. Work on uncertain-pose grasping with information
gathering \cite{hsiao2011}, conformal object-pose sets \cite{yang2023}, and reported failures from object
motion relative to a giver \cite{handover2025} already addresses nearby questions. They also noted robust
grasps across uncertain shapes \cite{lundell2019}, moving-object handover \cite{zhang2023}, and haptic
pose estimation under occlusion \cite{azulay2022}. These checks do not establish the proposed
matched-history effect for this pair, but they make the elementary certificate and sensing identity
insufficient as stand-alone contributions.

\section{Why the thread stopped}

The verifier accepted the conditional-risk and sensing calculations under their stated assumptions. It
paused the thread because P already states the basic matched-observation obstruction, while Q adds useful
ways to account for information, error and cost without supplying a tested handover result. The smallest
restart is to hold the donor wrist fixed, record the complete receiver-visible history, and externally
measure the marked object's pose through controlled in-hand rotations or reseating. Candidate matched
histories must then be checked for incompatible successful receiving actions. If they exist, held-out
receiving trials can compare a history-only or common-grasp choice with one new observation, including
its delay and cost. Until those measurements exist, the practical value of the connection is unknown.

\clearpage
\bibliographystyle{plainurl}
\bibliography{references}

\end{document}
